\documentclass[10pt,a4paper]{amsart}
\usepackage[margin=19mm]{geometry}
\usepackage{amsmath,amssymb,mathtools}
\usepackage[scr=rsfs,cal=euler]{mathalfa}
\usepackage{xfrac}
\usepackage[unicode,hidelinks,bookmarks=false]{hyperref}
\newcommand{\ie}{i.e.\ }
\newcommand{\cf}{cf.\ }
\newcommand{\resp}{respectively\ }
\newcommand{\Subsection}{\subsection}
\providecommand{\nobreakdash}{\nobreak\hbox{-}}
\setlength{\emergencystretch}{2em}
\multlinegap0pt
\usepackage{bm}
\usepackage{bbm}
\usepackage{dsfont}
\usepackage{xspace}
\usepackage{cancel}
\usepackage{mathtools}
\usepackage{amsthm}
\makeatletter
\@ifundefined{thm}{\newtheorem{thm}{Thm}}{}
\@ifundefined{lem}{\newtheorem{lem}[thm]{Lemma}}{}
\makeatother
\title[Some lower bounds for the maximal number of A-singularities ]{Some lower bounds for the maximal number of A-singularities in algebraic surfaces. II}
\author{Juan Garc\'ia Escudero}
\date{Source: arXiv:2604.15060v4 (CC BY 4.0)}
\begin{document}
\maketitle

\noindent\emph{Source and licence.} Some lower bounds for the maximal number of A-singularities in algebraic surfaces. II. Version arXiv:2604.15060v4 (CC BY 4.0), distributed under the Creative Commons Attribution 4.0 International licence, as recorded in the arXiv licence metadata for this exact version and shown on the abstract page https://arxiv.org/abs/2604.15060v4. Full rights notice: licenses/arxiv-2604.15060-CC-BY-4.0.txt.\par\smallskip

\noindent\textit{Editorial scope.} Counted unit: the Lem whose source label is \texttt{None} in the article (source0.tex lines 451--468, source label \texttt{None}), delivered together with the article's own complete original proof of it (source0.tex lines 469--595, one \texttt{proof} environment of 127 lines). No mathematics is written, repaired or re-derived: statement and proof are the article's own text line for line. Required context retained uncounted: none. Every definition, display and label of those lines is retained unchanged. Omitted from this package and not redistributed: 1--450, 596--800. Nothing inside a retained excerpt is omitted or altered: every retained body is delivered exactly as the article's lines are written, so no omission receipt of any kind accompanies this package. Citations: the retained excerpts carry no citation command at all, so no bibliography entry of the article is transcribed and no reference is redistributed. Reference adaptation: none was needed and none was made; every reference inside the delivered unit resolves inside it, so no unresolved reference or citation is produced. Printed slips kept literal: no printed word, symbol, hypothesis or number of the counted statement or of its proof was altered. Layout and class: the article's own class is \texttt{birkjour\_t2} with packages including ; the wrapper is the lane's pinned \texttt{amsart} editorial wrapper at 10pt on A4, which restates the theorem-like environments the retained text needs and adds the article's own macro definitions that the retained text needs (none), with extra wrapper packages (bm, bbm, dsfont, xspace, cancel, mathtools). The delivered numbering is the unit's own. Assets: no retained span contains a figure environment, a graphics-inclusion command, a PDF-page inclusion, a table or a TikZ picture; no image file is copied into this package, the package declares no asset dependency and no image file is redistributed with it. Licence: version arXiv:2604.15060v4 of this article is distributed under the Creative Commons Attribution 4.0 International licence, as recorded in the arXiv licence metadata for this exact version and shown on the abstract page https://arxiv.org/abs/2604.15060v4. A full rights notice is delivered at licenses/arxiv-2604.15060-CC-BY-4.0.txt. Characters: every retained excerpt is pure ASCII, so the delivered engine, XeLaTeX with its default Latin Modern text font, reports no missing character.
          \begin{lem} The series ${\mathcal{G}}^{(t)}_{d,\nu}(w), t=4,5,6$, can be extended to a larger series of polynomials, also denoted by ${\mathcal{G}}^{(t)}_{d,\nu}(w)$, having the following degrees $d=d(\nu, l, r, i)$
       \bigskip\par          
  (a)  $\nu= 3n+2$,  $n \in {\Bbb{Z}}^{\geq 1}$: $d(\nu, l, r, i)=(\nu+1)(l\nu + r+1) + 3i$,  $l \in {\Bbb{Z}}^{\geq 0}$,  $r=0,1,\dots 3n+1$, $i=0,1,\dots n$.
        \bigskip\par 
  (b) $\nu= 3n+3$,  $n \in {\Bbb{Z}}^{\geq 0}$:
  
  (1) $n=0$, $l \in {\Bbb{Z}}^{\geq 0}$, $r \in \{1,2,3\}$: $d(3, l, r, i)= 3(4l+r+1)+3i$, where $i=0$ if $r \in \{1,2\}$ and $i=0,1$, if $r=3$,
  
  (2) $n \in {\Bbb{Z}}^{\geq 1}$, $l \in {\Bbb{Z}}^{\geq 0}$, $r=0, \dots \nu$: $d(\nu, l, r, i)=\nu(l(\nu +1)+r+1)+ 3i$, with  $i=1,\dots n$ if $r=0$;  $i=0,1,\dots, n$ if $r \in \{1, \dots , \nu-1\}$ and $i=0,1,\dots, n+1$ if $r= \nu$.
      \bigskip\par 
  (c) $\nu= 3n+4$,  $n \in {\Bbb{Z}}^{\geq 0}$, $l_{k}=3k+b, b\in \{ 0, 1, 2 \}$ $k \in {\Bbb{Z}}^{\geq 0}$: 
  
  (1)  $l_{k}=0, r=-1$:  $d(\nu, l_{k}, r, i)= \nu +2+3i$, $i=0,1,\dots 2n+2$,
  
  (2) $l_{k}  \in {\Bbb{Z}}^{\geq 0}$, $r=0,1,\dots n$ if $b\in \{ 0, 1 \}$, $r=-1,0,\dots n$ if $b=2$: $d(\nu, l_{k}, r, i)=(\nu +1)(l_{k} (\nu-1)+3(r+\lfloor\frac{l_{k}+1}{3}\rfloor +1))+ 3i$, $i=0,1,\dots \nu$.
          
  The  new polynomials satisfy condition $(E)$.
  \end{lem}

           \begin{proof}     
               We denote  $C_{d}= \Big \lfloor \frac{d}{\nu+1} \Big \rfloor$, $D_{d}=\Big \lfloor \frac{d-1}{\nu} \Big \rfloor$, for a given value of $\nu$, and 
$U_{X}$, $X=C, D$ are the sets where $X_{d+3}- X_{d}=1$ if $i \in U_{X}$  and  $X_{d+3}- X_{d}=0$ if $i \notin U_{X}$. 

(a) For the cases  $\nu= 3n+2$ in Lemma 2.5 with $n \in {\Bbb{Z}}^{\geq 1}$ (the interpolation is not necessary when $n=0$) we consider $d(\nu, l, r, i)=(\nu+1)(l\nu + r+1) + 3i$, $l \in {\Bbb{Z}}^{\geq 0}$, $r=0,1,\dots 3n+1$, $i=0,1,\dots, n$, which satisfies $d(\nu, l+1, r, i)= d(\nu, l, r, i)+ \nu(\nu+1)$. The initial tree, corresponding to $l=r=i=0$, is represented in Fig.10(a). In this case $U_{C}=\{n\}, U_{D}=\{n-\Big \lfloor\frac{r+1}{3} \Big \rfloor \}$ if $r=0, 1, \dots \nu -2$ and $U_{D}=\{0, n\}$ if $r = \nu -1$. We also have $X_{d+3+\nu(\nu+1)}- X_{d+\nu(\nu+1)} = X_{d+3}- X_{d}, X=C, D$. 

 \bigskip\par
    
        (b)  For $\nu= 3n+3$, $n \in {\Bbb{Z}}^{\geq 0}$  the interpolated degrees $d(\nu, l, r, i)$ are
      \bigskip\par             
                 (1)  $n=0$, $l \in {\Bbb{Z}}^{\geq 0}$, $r \in \{1,2\}$: $d(3, l, r, i)= 3(4l+r+1)$,
                 
                 (2)  $n=0$, $l \in {\Bbb{Z}}^{\geq 0}$, $r=3$: $d(3, l, 3, i)= 3(4l+r+1)+3i$,  $i=0,1$, 
                 
                  (3) $n \in {\Bbb{Z}}^{\geq 1}$, $l \in {\Bbb{Z}}^{\geq 0}$, $r=0$: $d(\nu, l, 0, i)=\nu(l(\nu +1)+1)+ 3i$, $i=1,\dots n$,
                  
                   (4) $n \in {\Bbb{Z}}^{\geq 1}$, $l \in {\Bbb{Z}}^{\geq 0}$, $r \in \{1, \dots , \nu-1\}$: $d(\nu, l, r, i)=\nu(l(\nu +1)+r+1)+ 3i$, $i=0,1,\dots, n$, 
                   
                    (5) $n \in {\Bbb{Z}}^{\geq 1}$, $l \in {\Bbb{Z}}^{\geq 0}$, $r= \nu$: $d(\nu, l, \nu, i)=\nu(l+1)(\nu +1)+ 3i$, $i=0,1,\dots, n+1$, 
                   
                   
       \bigskip\par            
  \noindent               
 with the property $d(\nu, l+1, r, i)=d(\nu, l, r, i)+ \nu(\nu+1)$, $r \geq 0$.  The tree with the minimal degree $\nu+3$ ($d(3, 0, 1, 0)$ for $n=0$ and $d(\nu, 0, 0, 1)$ for  $n \geq 1$) is represented in Fig.10(b). The sets $U_{X}$ are 
      \bigskip\par 
    (1)  $n=0, l \in {\Bbb{Z}}^{\geq 0}$: 
    
    (1.1) $r=1, 2$: $U_{C}= U_{D}=\{ 0 \}$,
    
     (1.2) $r=3$: $U_{C}=\{ 1 \}, U_{D}=\{ 0,1 \}$,
    
     (2)  $n \in {\Bbb{Z}}^{\geq 1}, l \in {\Bbb{Z}}^{\geq 0}$: 
     
     (2.1) $r=0$: $U_{C}= U_{D}=\emptyset$,
    
 (2.2) $r=1, 2$: $U_{C}= U_{D}=\{0\}$, 
 
 (2.3) $r=3, \dots \nu-1$: $U_{C}=\{ \lfloor \frac{r}{3}  \rfloor\}$,  $U_{D}=\{0\}$,
 
 (2.4) $r=\nu$: $U_{C}=\{ n+1 \}$,  $U_{D}=\{0, n+1\}$.
    \bigskip\par 
In this case $X_{d+3+ \nu(\nu+1)}- X_{d+ \nu(\nu+1)} = X_{d+3}- X_{d}, X=C, D$. 

  \bigskip\par
(c) For $\nu= 3n+4$ the interpolated degrees $d(\nu, l_{k}, r, i)$, $n \in {\Bbb{Z}}^{\geq 0}$, $l_{k}=3k+b, b\in \{ 0, 1, 2 \}$, $k \in {\Bbb{Z}}^{\geq 0}$ are 
      \bigskip\par 
  (1) $l_{k}=0$, $r=-1$: $d(\nu, 0, -1, i)=\nu+2+3i$,  $i=0,\dots 2n+2$,
  
   (2) $l_{k}=3k+b$, $b\in \{ 0, 1 \}$: $d(\nu, l_{k}, r, i)=(\nu +1)(l_{k} (\nu-1)+3(r+\lfloor\frac{l_{k}+1}{3}\rfloor +1))+ 3i$, $r = 0, \dots, n$, $i=0,\dots \nu$,
   
     (3) $l_{k}=3k+b$, $b=2$: $d(\nu, l_{k}, r, i)=(\nu +1)(l_{k} (\nu-1)+3(r+\lfloor\frac{l_{k}+1}{3}\rfloor +1))+ 3i$, $r = -1, 0, \dots, n$, $i=0,\dots \nu$,

       \bigskip\par 
  \noindent            
 with  $d(\nu, l_{k+1}, r, i)=d(\nu, l_{k}, r, i)+ 3\nu(\nu+1)$ and $X_{d+3+3 \nu(\nu+1)}- X_{d+3 \nu(\nu+1)} = X_{d+3}- X_{d}, X=C, D$, when $d \geq 3(\nu+1)$. The initial tree, with $d(\nu, 0, -1, 0)=\nu+2$, can be seen in Fig.10(c). In this case the sets $U_{X}$ are  
       \bigskip\par 
 (1) $n \in {\Bbb{Z}}^{\geq 0}$, $ l_{k}=0, r=-1$:   $U_{C}=\{n+1, 2n+2\}$,  $U_{D}=\{n, 2n+2\}$ 
 
 (2) $n=0$,  $l_{k}=3k+b, b\in \{ 0, 1, 2 \}$, $k \in {\Bbb{Z}}^{\geq 0}$: 
 
  (2.1) $l_{k}=3k+b, b\in \{ 0, 1\}$, $r=0$: $U_{C}=\{1, 3, 4\}$, $U_{D}=\{0, b+1, 3, 4\}$ 
  
   (2.2) $l_{k}=3k+2$: $U_{C}=\{1, 3, 4\}$ and $U_{D}=\{1, 2, 3\}$ if $r = -1$ and $U_{D}=\{0, 1, 2, 4\}$ if $r = 0$
 
   (3) $n \in {\Bbb{Z}}^{\geq 1}$, $l_{k}=3k+b, b\in \{ 0, 1, 2 \}$, $k \in {\Bbb{Z}}^{\geq 0}$: $U_{C}=\{n+1, 2n+3, 3n+4\}$ and 
 
 (3.1) $l_{k}=3k+b, b\in \{ 0, 1\}$: $U_{D}=\{n-r, 2n-r+b+1, 3n-r+3\}$ if $r = 0, \dots n -1$ and 
 
 $U_{D}=\{0, n+b+1, 2n+3, 3n+4\}$ if $r = n$,
 
 (3.2) $l_{k}=3k+2$:  $U_{D}=\{n-r, 2n-r+1, 3n-r+2\}$ if $r = -1, 0, \dots n -1$ and 
 
 $U_{D}=\{0, n+1, 2n+2, 3n+4\}$ if $r = n$.
 


    \bigskip\par 
 We now apply a series of transformations $w= \alpha$, $\beta$, $\bar{\alpha}$, $\bar{\beta}$ to $\mathcal{T}_{\nu}(d, N_{-1}, N_{1})$ giving $\mathcal{T}_{\nu}(d+3, N_{-1}+a, N_{1}+b), a,b \in \{-1,0,1\}$. The transformation $\alpha$ consists in adding black vertices to a white one in such a way that it is transformed into one with $\nu+1$ edges, hence $a=0$, $b=1$. With $\beta$ a black vertex is converted into one with $\nu+1$ adjacent edges in such a way that the resulting tree is $\mathcal{T}_{\nu}(d+3, N_{-1}+1, N_{1}-1)$. The transformation $\bar{\alpha}$ is like $\alpha$ but the transformed vertex has less than $\nu+1$ adjacent edges ($a=0$, $b=0$) and  $\bar{\beta}$ is like $\beta$ but without subtracting edges from a white vertex ($a=1$, $b=0$). By employing the notation  $Y= (C_{d+3}- C_{d}, D_{d+3}- D_{d}; w)$, which means that $w$ is applied when the first two terms of the triple take the indicated values, we have $Y \in$ $\{(0,1; \alpha)$,  $(1,0; \beta)$, $(0, 0; \bar{\alpha})$,  $(1, 1; \bar{\beta})\}$. We analyse the words representing the sequence of trees, which can be obtained from the sets $U_{X}$ described above.
 
      \bigskip\par 
 
(a) $\nu= 3n+2$: If we start with the tree with $n=1, l=1$ in Lemma 2.5 then a sequence of allowed words is $W^{c}$, $W^{c}\bar{\alpha}$, $W^{c}\bar{\alpha}\bar{\beta}$, $\dots$, $W^{c+1}$, $W^{c+1}\bar{\alpha}$, $\dots $, $c \in  {\Bbb{Z}}^{\geq 1}$, where $W=(\bar{\alpha}\bar{\beta})^{2}(\alpha\beta)^{2}\alpha\bar{\beta}$ describes the chain of transformations from $d=36$ to $d=66$. The words appear repeated in cycles with one cycle from $l$ to $l+1$. In this Lemma we also incorporate $l=0$, which corresponds here to $d=6$, and then the word representing the transformations to the tree with $d=36$ is also $W$. In what follows $W$ represents the complete word within a cycle, and the allowed words are all its prefixes. If we use the product notation for concatenation, then for $n \geq 2$  the words within a cycle are the prefixes of 

  \begin{equation} 
 W=\bar{Q}_{n}^{2}\left(\prod^{n-1} _{p=1}P_{n, p-1}^{3}\right) (\alpha Q_{n-1})^{2} \alpha \bar{Q}_{n-1}
           \end{equation} 
with $P_{n,p}=\bar{\alpha}^{n-p-1}\alpha \bar{\alpha}^{p} \beta, Q_{n}= \bar{\alpha}^{n}\beta, \bar{Q}_{n}= \bar{\alpha}^{n}\bar{\beta}$.
     \bigskip\par 
 
(b) $\nu= 3n+3$: The cycle in this case is from $l$ to $l+1$ and its word $W$ is 
      \bigskip\par 
(1) $n =0$, $l \in  {\Bbb{Z}}^{\geq 0}$

      \begin{equation} 
W=  \bar{Q}_{0}^{2}  \alpha \bar{Q}_{0}
          \end{equation}  

(2) $n \in  {\Bbb{Z}}^{\geq 1}$, $l \in  {\Bbb{Z}}^{\geq 0}$
      \begin{equation} 
W=  \bar{Q}_{n}^{2} \bar\alpha^{n} \left( \prod^{n-1} _{p=0}( \alpha {Q}_{p } \bar\alpha^{n-p -1})^{3}\right) \alpha  \bar{Q}_{n}
          \end{equation}  

 The possible words are $W^{c}$, $c \in  {\Bbb{Z}}^{\geq 1}$ and its prefixes.
      \bigskip\par 

(c) $\nu= 3n+4$:  for $l_{k}=3k+b, b\in \{ 0, 1, 2 \}$, the  words  are  $P_{n+1,0} \bar{Q}_{n} \prod _{k \geq 0} W_{3k}W_{3k+1}W_{3k+2}$, and its prefixes. The cycle occurs from $l_{k}$ to $l_{k+1}$ which corresponds to  $W=W_{3k}W_{3k+1}W_{3k+2}$.
      \bigskip\par 
(1) $n =0$, $l \in  {\Bbb{Z}}^{\geq 0}$

   \begin{equation} 
  W_{3k}=\alpha\bar{\beta}\bar{\alpha}\bar{\beta}^2, W_{3k+1}=\alpha\beta\alpha\bar{\beta}^2, W_{3k+2}=\bar{\alpha}\bar{\beta}\alpha\bar{\beta}\beta\alpha\bar{\beta}\alpha\beta\bar{\beta}
           \end{equation} 

(2) $n \in  {\Bbb{Z}}^{\geq 1}$, $l \in  {\Bbb{Z}}^{\geq 0}$
    \begin{equation} 
  W_{3k}=\left(\prod^{n-1} _{p=0} \bar{\alpha}P_{n,p}P_{n+1,p+1}P_{n,p}\right)\alpha\bar{Q}_{n}\bar{Q}_{n+1} \bar{Q}_{n}  
           \end{equation} 
  \begin{equation} 
W_{3k+1}=\left(\prod^{n-1} _{p=0} \bar{\alpha}P_{n,p}P_{n+1,p}P_{n,p}\right)\alpha Q_{n}  \alpha  \bar{Q}_{n}^{2}
         \end{equation}  
  
      \begin{equation} 
    W_{3k+2}=\bar{\alpha}  \bar{Q}_{n} P_{n+1,0}P_{n,0}\left(\prod^{n-2} _{p=0} \bar{\alpha}P_{n,p}P_{n+1,p+1}P_{n,p+1}\right)\bar{\alpha} \alpha Q_{n-1} \alpha 
    \bar{Q}_{n}Q_{n}\alpha\bar{Q}_{n}\bar{\alpha}^{n} \alpha \beta \bar{Q}_{n}.
             \end{equation}  
 
  \end{proof}

\end{document}
